\chapter{Il richiamo delle Ande, e i suoi limiti}
\label{cap:qhapaq_nan}

Questo libro è partito dalla Región de Atacama, e alle Ande conviene tornare adesso, non più come fonte di un contrasto visivo, ma come laboratorio di idee. Le popolazioni andine hanno lasciato, oltre alle incisioni rupestri da cui si è aperto questo racconto, un'altra eredità straordinaria: una rete di percorsi che attraversa l'intero continente sudamericano lungo la dorsale delle Ande, conosciuta come Qhapaq Ñan, il grande cammino inca. È naturale chiedersi se questo caso, così ben studiato e così spettacolare, possa insegnare qualcosa di utile per il problema di questo libro: come ricostruire i percorsi che i cacciatori preistorici seguivano sulle Alpi venete.

La risposta, come si vedrà, è più sfumata di un semplice sì. Il Qhapaq Ñan insegna qualcosa di importante, ma non nel modo in cui ci si potrebbe aspettare a prima vista.

\section{Le dimensioni del problema}

Prima di ragionare sulla logica del tracciato, vale la pena avere un'idea delle sue dimensioni reali. Il dossier con cui il Qhapaq Ñan è stato iscritto nella lista del patrimonio mondiale dell'UNESCO nel 2014 documenta 137 aree componenti e 308 siti archeologici associati, per un totale di 616,06 chilometri di percorso formalmente censito \citep{unesco_qhapaqnan}: una cifra che riguarda solo la parte della rete formalmente riconosciuta e protetta, non l'intera estensione della rete stradale inca, stimata su scala continentale in decine di migliaia di chilometri. Lo stesso dossier descrive il cammino come un'opera che collega le vette innevate delle Ande, a oltre seimilaseicento metri di quota, fino alla costa, attraversando foreste pluviali calde, valli fertili e deserti assoluti: un'unica infrastruttura capace di funzionare in condizioni ambientali quasi agli estremi opposti.

\section{Un percorso che non segue solo la fisica}

Il primo istinto, di fronte al problema di ricostruire un percorso antico, è pensare in termini di efficienza: un cammino tende a seguire la via più economica, quella che minimizza lo sforzo fisico di chi lo percorre. Esistono, in effetti, formule matematiche pensate proprio per stimare questo sforzo in funzione della pendenza del terreno, sviluppate originariamente per escursionisti moderni e poi adottate largamente in archeologia per ricostruire percorsi antichi.

Il Qhapaq Ñan mette in crisi, quasi subito, questa idea. Lo stesso dossier UNESCO descrive la rete come un'opera che dimostra maestria nel risolvere gli innumerevoli problemi posti dalla variabilità del paesaggio andino, per mezzo di tecnologie costruttive diversificate: ponti, scalinate, fossati di drenaggio, pavimentazioni in ciottoli. In molti tratti il percorso sale con pendenze che un semplice calcolo di minimizzazione dello sforzo avrebbe scartato, sostenute da queste stesse opere di ingegneria, che rendevano percorribile ciò che altrimenti non lo sarebbe stato. Diverse fonti tecniche concordano su un punto specifico: nei tratti più ripidi, gli ingegneri inca costruivano scalinate intagliate nella roccia, oppure facevano deviare il tracciato a zig-zag, invece di cercare la linea a pendenza più dolce come farebbe un modello di costo energetico puro.

Il Qhapaq Ñan, insomma, non è il prodotto di migliaia di scelte individuali che, sommate, convergono verso il percorso più comodo. È, in larga misura, un'infrastruttura pianificata da uno stato, con finalità che vanno oltre la sola economia dello spostamento: amministrazione, comunicazione militare, prestigio, religione, e un controllo del territorio che si manifesta anche nella standardizzazione delle tecniche costruttive imposta dallo stato inca su regioni e contesti locali molto diversi tra loro.

\section{Una controprova empirica, non solo un'osservazione}

Chi sospetta che i percorsi umani non seguano sempre il tragitto energeticamente più conveniente ha a disposizione, oltre al caso monumentale del Qhapaq Ñan, anche una controprova più piccola ma metodologicamente istruttiva. Uno studio ormai classico di Marcos Llobera e Timothy Sluckin ha applicato l'algoritmo di Dijkstra, lo stesso usato per calcolare percorsi a costo minimo su un modello digitale del terreno, a un'area montuosa del Galles, confrontando i tracciati calcolati con i sentieri pedonali effettivamente esistenti sul territorio \citep{llobera2004}.

Il risultato è che i sentieri reali non seguono il tempo minimo di percorrenza, come un primo istinto potrebbe far pensare, ma un compromesso diverso, legato non al tempo ma al costo metabolico dello sforzo fisico: un percorso che minimizza il tempo tende a salire più ripidamente possibile appena il vantaggio in velocità supera lo svantaggio in pendenza, mentre un percorso reale tende a preferire tragitti più lunghi ma meno faticosi. Lo stesso studio mostra che il parametro che meglio descrive questo compromesso varia da sentiero a sentiero, segno che nemmeno tra escursionisti moderni, su un terreno privo di vincoli amministrativi o simbolici, esiste un'unica funzione di costo universale.

Questo studio, condotto su un caso molto più modesto del Qhapaq Ñan, rafforza comunque lo stesso punto metodologico: anche quando l'unico obiettivo è muoversi da un punto all'altro, senza le complicazioni di un'infrastruttura di stato, il comportamento umano reale si allontana dal puro calcolo del tempo minimo. A maggior ragione, un'infrastruttura come il Qhapaq Ñan, costruita per ragioni che vanno ben oltre lo spostamento individuale, non può essere ridotta a un semplice esercizio di minimizzazione energetica.

\section{Perché questo limite non si applica a tutti i casi}

A prima vista, queste osservazioni potrebbero sembrare un argomento contro l'uso di qualunque modello di costo energetico per ricostruire percorsi antichi, comprese le Alpi venete preistoriche. Sarebbe però una conclusione troppo affrettata, perché il Qhapaq Ñan e i cacciatori-raccoglitori delle Alpi venete rappresentano due situazioni profondamente diverse, per almeno due ragioni.

La prima ragione è di scala organizzativa. Il grande cammino inca è opera di uno stato con capacità di mobilitare manodopera su larga scala per costruire infrastrutture permanenti, come mostrano le centinaia di siti associati e le centinaia di chilometri di opere murarie censite dall'UNESCO. I gruppi di cacciatori-raccoglitori che frequentavano le Alpi venete in epoca preistorica, come si è visto nei capitoli precedenti, erano piccoli nuclei familiari allargati, senza le risorse né l'organizzazione necessarie per costruire gradini nella roccia o sistemi di drenaggio. I loro spostamenti, con ogni probabilità, seguivano vie naturali già esistenti, non infrastrutture create appositamente.

La seconda ragione, più sottile, riguarda la natura di ciò che si vuole ricostruire. Il Qhapaq Ñan è un percorso umano, costruito da esseri umani per finalità umane, comprese quelle simboliche e amministrative che nessun modello fisico può catturare. Anche lo studio più modesto sui sentieri gallesi conferma che persino un semplice sentiero pedonale moderno si discosta dal puro calcolo energetico. Ma nel problema che interessa questo libro, prima ancora dei percorsi degli uomini, contano i percorsi degli animali che gli uomini cacciavano. E un animale, a differenza di un imperatore inca che pianifica una strada per collegare il proprio impero, e a differenza di un escursionista che sceglie un sentiero anche per ragioni estetiche o di abitudine, non ha ragioni simboliche o amministrative per scegliere un tragitto. Il costo energetico è uno dei principali vincoli che ne guidano il movimento, insieme alla qualità dell'habitat, alla copertura nevosa, alla disponibilità di cibo, alle esigenze riproduttive e al rischio di predazione: variabili che agiscono tutte nella stessa direzione di minimizzare lo sforzo e massimizzare la resa, e che un modello di pendenza descrive nella sua componente più elementare.

\section{Cosa si porta a casa da questo confronto}

Il Qhapaq Ñan, allora, non offre un modello da copiare direttamente, ma un avvertimento prezioso, rafforzato da una controprova empirica indipendente, da tenere bene a mente per il resto di questa parte del libro: un modello di costo energetico può descrivere bene il movimento di un animale, ma diventa rischioso quando si applica a un percorso costruito o scelto dall'uomo per ragioni che vanno oltre la sola economia dello sforzo fisico. Prima di applicare un simile modello, occorre sempre chiedersi: sto ricostruendo il movimento di un animale, guidato dalla sola fisiologia, o il movimento di un essere umano, che porta con sé anche significati, convenzioni, scelte non riducibili alla sola pendenza del terreno?

È una distinzione che sembra ovvia una volta enunciata, ma che è facile perdere di vista quando ci si getta a capofitto nella costruzione di un modello. Il prossimo capitolo riprende proprio da qui, per chiedersi con più precisione quando un modello di costo energetico sia davvero applicabile, e a quale delle due categorie, animale o umana, appartenga il problema specifico di questo libro.
